\section{Introducció}
Text de prova\cite{Prova}

Per posar signes matematics a les sections o subsections:
\section{Some title with \texorpdfstring{$\mu$}{μ} and \texorpdfstring{$\sigma$}{σ}}

Command per a compilar:
\begin{verbatim}
    latexmk main.tex
\end{verbatim}


Command per a que es compili sol cada cop que guardo un fitxer:
\begin{verbatim}
    latexmk -pvc main.tex
\end{verbatim}

\section{Per escriure codi amb minted}
Minted et fa el syntax highlighting

\begin{minted}{python}
import numpy as np

for i in range(10):
    pass

a = np.sum()

\end{minted}

\begin{minted}{C++}
#include <iostream>
#include <raylib.h>

#define WIDTH 900

void init(int h){
    InitWindow(WIDTH, h, "Window Title");
}

int main(){std::cout << "HW!\n";}

\end{minted}


Per fer dues figures en el lloc d'una es fa:
\begin{figure}[H]
    \centering

    \begin{subfigure}[b]{0.48\textwidth}
        \centering
        \includegraphics[width=\textwidth]{path_to_figure}
        \caption{one figure}
        \label{fig:first}
    \end{subfigure}
    \hfill
    \begin{subfigure}[b]{0.48\textwidth}
        \centering
        \includegraphics[width=\textwidth]{path_to_another_figure}
        \caption{Another figure}
        \label{fig:second}
    \end{subfigure}

    \caption{Comparison of the two figures.}
    \label{fig:combined}
\end{figure}

I si es volen verticals es fa: (simplement posant que ocupin la totalitat de l'amplada de la pàgina amb \\textwidth)
\begin{figure}[H]
    \centering

    \begin{subfigure}{\textwidth}
        \centering
        \includegraphics[width=0.95\textwidth]{path_to_one_figure}
        \caption{First figure}
    \end{subfigure}

    \vspace{0.5cm}

    \begin{subfigure}{\textwidth}
        \centering
        \includegraphics[width=0.95\textwidth]{path_to_another_figure}
        \caption{Second figure}
    \end{subfigure}

    \caption{Two vertically stacked figures.}
\end{figure}
